\chapter[\emj{🚰}~Flujo Máximo (Ford-Fulkerson / Edmonds-Karp)]{\emj{🚰}~Flujo Máximo (Ford-Fulkerson / Edmonds-Karp)}

\begin{dsaquees}

Una red de tuberías 🚰 desde una fuente (source) hasta un desagüe
(sink). Cada tubo aguanta un caudal máximo (capacidad). Pregunta: ¿cuánta
agua total puedes empujar de la fuente al desagüe al mismo tiempo?

La idea: mientras exista un camino con espacio libre de la fuente al
desagüe, manda tanta agua como el tubo más estrecho de ese camino
permita. Repite hasta que ya no quede ningún camino con espacio. La
suma de todo lo que enviaste es el flujo máximo.

El truco genial: dejas "aristas de retorno" que permiten DESHACER
decisiones previas si conviene reencaminar el agua.

\end{dsaquees}

\begin{dsaotraforma}

El problema de \textbf{flujo máximo} busca enviar la mayor cantidad de
flujo posible desde un nodo fuente $s$ a un sumidero $t$ en un grafo
dirigido con capacidades en las aristas, sin exceder ninguna capacidad.

\textbf{Ford-Fulkerson} repite: encuentra un "camino de aumento"
(augmenting path) de $s$ a $t$ con capacidad residual $>0$, empuja por él
el mínimo de sus capacidades (el cuello de botella) y actualiza el grafo
residual (resta en la arista, suma en su arista inversa). Termina cuando
no hay más caminos.

\textbf{Edmonds-Karp} es Ford-Fulkerson eligiendo SIEMPRE el camino de
aumento más corto vía BFS, lo que garantiza $O(V E^2)$.

Teorema \textbf{max-flow min-cut}: el flujo máximo iguala la capacidad del
corte mínimo que separa $s$ de $t$.

\end{dsaotraforma}

\begin{dsadiagrama}
Red (capacidades):
        10        10
   S ───────> A ───────> T
   │          ^          ^
   │10        │ 5        │10
   v          │          │
   B ─────────┘          │
        (B->A cap 5, B->T cap 10 ...)

Camino de aumento S->A->T: cuello de botella min(10,10)=10 -> envía 10
Camino de aumento S->B->T: cuello min(10,10)=10 -> envía 10
No hay más caminos con espacio -> Flujo máximo = 20

Grafo residual: cada envío resta capacidad y crea arista inversa
(para poder "devolver" flujo si un camino mejor lo requiere).
\end{dsadiagrama}

\begin{dsacomofunciona}
\begin{verbatim}
GRAFO RESIDUAL evolucionando (flujo/capacidad en cada arista)

Red:   S →10→ A →10→ T        S →10→ B →10→ T

Inicio (residual = capacidad):
   S=10=>A   A=10=>T   S=10=>B   B=10=>T      flujo total = 0

Camino de aumento 1:  S→A→T, cuello = min(10,10) = 10
   S: 0/10 →A  (residual 0, inversa A→S = 10)
   A: 0/10 →T
   flujo total = 10

Camino de aumento 2:  S→B→T, cuello = min(10,10) = 10
   S: 0/10 →B
   B: 0/10 →T
   flujo total = 20

BFS ya no encuentra camino con residual > 0  →  FIN
Flujo máximo = 20

MIN-CUT equivalente:  aristas saturadas que separan S de T
   {S→A, S→B} capacidad 10+10 = 20  =  flujo máximo  OK (max-flow=min-cut)
\end{verbatim}
\end{dsacomofunciona}

\begin{lstlisting}
EDMONDS-KARP (BFS para el camino de aumento)
1. Matriz de capacidad residual cap[u][v].
2. BFS de s a t por aristas con cap > 0, guardando el padre.
3. Si llegas a t: el cuello de botella es el mínimo cap del camino.
   Resta ese flujo en las aristas y súmalo en las inversas.
4. Repite hasta que el BFS ya no alcance t. Suma total = flujo máximo.

PSEUDOCÓDIGO
──────────────────────────────────────────────────────────────

función maxFlow(cap, s, t):
    flujo = 0
    REPETIR:
        padre = BFS(cap, s, t)          // camino con capacidad > 0
        SI t no alcanzable: romper
        cuello = min de cap a lo largo del camino s..t
        PARA cada arista (u,v) del camino:
            cap[u][v] -= cuello
            cap[v][u] += cuello         // arista de retorno
        flujo += cuello
    return flujo

CÓDIGO C++ (Edmonds-Karp)
──────────────────────────────────────────────────────────────

#include <vector>
#include <queue>
#include <climits>
#include <cstring>
using namespace std;

int V;
long long bfs(vector<vector<long long>>& cap, int s, int t,
              vector<int>& parent) {
    fill(parent.begin(), parent.end(), -1);
    parent[s] = s;
    queue<pair<int,long long>> q;        // {nodo, flujo hasta aquí}
    q.push({s, LLONG_MAX});
    while (!q.empty()) {
        auto [u, f] = q.front(); q.pop();
        for (int v = 0; v < V; v++) {
            if (parent[v] == -1 && cap[u][v] > 0) {
                parent[v] = u;
                long long nf = min(f, cap[u][v]);
                if (v == t) return nf;    // llegó al sumidero
                q.push({v, nf});
            }
        }
    }
    return 0;                             // no hay camino
}

long long maxFlow(vector<vector<long long>>& cap, int s, int t) {
    long long flow = 0, nf;
    vector<int> parent(V);
    while ((nf = bfs(cap, s, t, parent)) > 0) {
        flow += nf;
        int v = t;
        while (v != s) {                  // recorre el camino al revés
            int u = parent[v];
            cap[u][v] -= nf;              // consume capacidad
            cap[v][u] += nf;              // arista de retorno
            v = u;
        }
    }
    return flow;
}

COMPLEJIDAD (Big-O)
──────────────────────────────────────────────────────────────

  Algoritmo         │ Tiempo       │ Nota
  ──────────────────┼──────────────┼──────────────────────
  Edmonds-Karp      │ O(V · E^2)   │ BFS, capacidades enteras
  Dinic             │ O(V^2 · E)   │ mucho más rápido en la práctica

* Ford-Fulkerson genérico puede no terminar con capacidades irracionales.
\end{lstlisting}

\begin{dsaentrevistas}

✅ Max-flow min-cut: el flujo máximo = capacidad del corte mínimo. Muchos
   problemas "¿mínimo para desconectar/separar?" se reducen a max-flow.
   Sabértelo abre la puerta a modelar problemas difíciles.

💡 Matching bipartito máximo se resuelve con max-flow: fuente -> lado A
   -> lado B -> sumidero, capacidades 1. El flujo máximo es el número de
   emparejamientos.

⚠️ Aristas de retorno son OBLIGATORIAS: sin la capacidad inversa, el
   algoritmo no puede "arrepentirse" y da respuestas menores a la óptima.

🔑 Elige BFS (Edmonds-Karp): usar DFS sin control puede degradar el
   tiempo enormemente. BFS asegura caminos más cortos y el límite
   $O(VE^2)$.

\end{dsaentrevistas}

\begin{dsaresumen}

• Flujo máximo: máximo caudal de la fuente s al sumidero t sin pasarse.
• Ford-Fulkerson: repite empujar flujo por caminos de aumento.
• Edmonds-Karp: elige el camino más corto con BFS; O(V·E²).
• Grafo residual con aristas de retorno para reencaminar flujo.
• Teorema max-flow min-cut: flujo máx = corte mínimo.
• Modela matching bipartito y problemas de separación mínima.
════════════════════════════════════════════════════════════════

\end{dsaresumen}
